\section{Выводы}

В данной работе были экспериментально исследованы волоконно-оптические датчики деформации на основе SMS-структур с различными длинами многомодового участка. Основной целью было определение их спектральных характеристик под воздействием продольной нагрузки и анализ чувствительности.

Были экспериментально определены следующие ключевые характеристики:

\begin{itemize}
    \item Чувствительность датчиков к продольному растяжению (в пересчете на относительное удлинение) варьируется в диапазоне от $-15.7$ до $-29.2$ нм/\% для датчика с $L_{MMF} = 4.97$ см, от $-20.2$ до $-24.3$ нм/\% для $L_{MMF} = 10.2$ см и от $-16.1$ до $-18.3$ нм/\% для $L_{MMF} = 20.5$ см.
    \item Оптимальная длина многомодовой секции для достижения стабильной и воспроизводимой чувствительности составляет порядка $10$ см.
    \item Стандартная чувствительность волоконных датчиков деформации на основе Брэгговских решеток (FBG) составляет порядка 12 нм/\%, что несколько меньше чувствительности исследованных SMS-датчиков.
    \item Количество провалов в спектре пропускания увеличивается с ростом длины многомодовой секции SMS-датчика.
    \item Визуальная оценка числа мод в одномодовом волокне (d=9 мкм) составила порядка 3 мод, в то время как для многомодового волокна (d=62.5 мкм) наблюдалось значительно большее число интерферирующих мод (много больше 10).
    \item Экспериментально определенное угловое расхождение спекл-картины для одномодового волокна (d=9 мкм) составило 0.171 рад, что сравнимо с расчетным размером пятна Эйри (0.176 рад). Для многомодового волокна (d=62.5 мкм) угловое расхождение (0.318 рад) значительно превышает размер пятна Эйри (0.025 рад).
\end{itemize}

В результате подтверждена линейная зависимость сдвига спектральных минимумов от приложенной нагрузки, что является важным свойством для практического применения датчиков. Определенная чувствительность SMS-датчиков к продольному растяжению достигает высоких значений, превосходящих стандартную чувствительность датчиков на основе Брэгговских решеток.

Анализ спекл-картин наглядно продемонстрировал фундаментальное различие в распространении света в одномодовом (маломодовом) и многомодовом волокнах. Наблюдаемое в многомодовом волокне большое количество интерферирующих мод и сильная межмодовая интерференция, подтвержденная сравнением углового расхождения с пятном Эйри, приводит к дисперсии и уширению оптических импульсов. Этот эффект ограничивает дальность высокоскоростной передачи данных, обуславливая предпочтительное использование одномодовых волокон в магистральных линиях связи. Таким образом, работа позволила не только количественно охарактеризовать SMS-датчики деформации, но и наглядно проиллюстрировать различия в распространении света, определяющие их области применения.
